% Readability revision: the reviewed construction and domain are preserved.
% Local allocation residual r is renamed xi to distinguish the reuse term.
\subsection{Certifying a whole policy family}\label{sec:family-capacity}
We now optimize over policies. Here $\C$ is the fixed view; $w^*,\pi_0^*$ mark exact input probabilities, while $m^*,\pi^*$ will denote returned optimizers. The family $\Pi$
contains all randomized policies on the available actions. Its \emph{capacity}
is the largest worst-case improvement attainable by any policy.
Fix a nonempty compact convex view $\C\subseteq[0,1]^{KL}$,
exact probability weights $w^*\ge0$, $\sum_xw_x^*=1$, and reference
$\pi_0^*\in\Pi$. Define
\begin{align}
 \mathcal R(\C)&=\max_{\pi\in\Pi}\min_{m\in\C}
   \sum_{x,a}w_x^*(\pi_{xa}-\pi_{0,xa}^*)m_{xa},\label{eq:capacity}\\
 F(m)&=\sum_xw_x^*\left[\max_a m_{xa}
             -\sum_a\pi_{0,xa}^*m_{xa}\right].\label{eq:family-upper}
\end{align}
The inner minimum in~\eqref{eq:capacity} gives one policy's worst allowed case; the outer maximum chooses the best such policy. The reference makes $\mathcal R\ge0$. At one allowed $m$, $F(m)$ chooses the best action per context.
Every checked $m\in\C$ gives $\mathcal R(\C)\le F(m)$: each policy's
minimum is at most its payoff at this point, whose policy maximum is $F(m)$.
Thus $U(c)$ evaluates one policy, while $F(m)$ maximizes over policies at that point. The latter bounds the best robust guarantee; $F(m)\le\delta$ excludes all guarantees above $\delta$.
The witness need not equal the true means; it does not bound actual effects.

\begin{proposition}[Zero capacity in a fixed view]\label{prop:zero-capacity}
$\mathcal R(\C)=\min_{m\in\C}F(m)$. If all $w_x^*>0$ and
$\pi_{0,xa}^*>0$, then $\mathcal R(\C)=0$ exactly when $\C$ contains
a vector constant across actions within every context.
\end{proposition}
\noindent\textit{Proof.} Compact convex domains and a continuous bilinear payoff
permit minimax exchange~\cite{Sion1958Minimax}; row maximization gives $F$.
Its nonnegative summands, under strict positivity, vanish exactly at constant
rows. Compactness attains the minimum.\hspace*{\fill}\nobreakspace$\square$

Without strict positivity, constant points remain sufficient for zero capacity,
but the converse need not hold.

For nonnegative single-context moment rows, even overlapping ones, form
\[
 I_x=[\max_a\ell_{xa},\min_a u_{xa}]
       \cap\!\bigcap_{j:x_j=x}[a_j/\kappa_j,b_j/\kappa_j].
\]
Here $\kappa_j=\sum_iT_{ji}>0$; zero rows require $a_j\le0\le b_j$.
Choose one $v_x\in I_x$ per context, set every $m_{xa}=v_x$, and check the
original constraints. This \emph{constant-action witness} proves zero capacity
without optimizing. Under the proposition's positivity conditions,
an empty intersection proves $\mathcal R>0$, not $\mathcal R>\delta$ for
positive $\delta$.

Partial evidence can bracket the family without full optimization. The maximum
of checked lower bounds for any legal policies, together with the reference's
zero, is $L_{\rm saved}\le\mathcal R$. A nonempty set $M_{\rm saved}$ of checked
points gives $U_{\rm saved}=\min_{m\in M_{\rm saved}}F(m)\ge\mathcal R$.
Share them only for identical $\C,w^*,\pi_0^*$.
Then $L_{\rm saved}>\delta$ certifies existence, $U_{\rm saved}\le\delta$ excludes it,
and otherwise the family remains unresolved. Equation~\eqref{eq:recovery-decomposition}
separates certificate slack from restricted policy choice.

\subsection{Exact capacity for partitioned moments}\label{sec:partition-capacity}
For one context, $F(m)$ is the largest action mean minus the reference payoff. Minimax therefore asks us to minimize $t-p^\top m$, with $m_i\le t$ and reference probabilities $p=\pi_{0,x}^*$. Our construction has two tasks: find the smallest family upper bound, then recover a policy attaining it. For disjoint groups, the first task reduces to capping means and allocating below that cap. The second recovers a proof using the \emph{original} constraints, so the auxiliary cap cannot silently narrow the guarantee. Minimax and continuous allocation are classical ingredients~\cite{Eiselt1986ContinuousGLB,Kuno1991ContinuousKnapsack}.

\begin{theorem}[An attaining policy with an original-view certificate]\label{thm:partition-certificate}
Let $\C$ be a nonempty rational box-and-moment view with
$0\le\ell\le u\le1$. Suppose every nonzero moment row is nonnegative,
belongs to one context, and has positive support disjoint from every other
row. For any rational probability weights $w^*$ and reference policy $\pi_0^*$,
including zero entries, the three-step construction returns
$m^*\in\C$, $\pi^*\in\Pi$ and nonnegative original-row multipliers
$\eta_L,\eta_U$ such that, with $c^*_{xa}=w_x^*(\pi^*_{xa}-\pi^*_{0,xa})$,
\[
 D(c^*;\eta)=c^{*\top}m^*=F(m^*)=\mathcal R(\C).
\]
The dense implementation uses $O(dJ+d^2)$ exact arithmetic/comparison
operations; the independent receiver uses $O(dJ+d)$. These are unit-cost
rational-operation counts, excluding coefficient-dependent bit costs,
encoding and hashing.
\end{theorem}

The certificate closes the gap between a policy we can achieve and a family bound no policy can exceed:
$D(c^*;\eta)\le\min_{m\in\C}c^{*\top}m\le\mathcal R\le F(m^*)$.
Equal checked endpoints prove optimality independently of how the producer found them. Check the domain before construction; reject unsupported rows.
Zero rows require $a_j\le0\le b_j$ and receive zero multipliers.
Ties use index order. The following three steps establish the theorem.

\textbf{Step 1: Find the smallest feasible cap.}

Within this algorithm, $p=\pi_{0,x}^*$ is the reference distribution, distinct
from observed success $p_{xa}$. Write $\ell\le m\le u$ and
$a_j\le T_jm\le b_j$. Set
\[
 \begin{aligned}
 t_\ell&=\max_i\ell_i,\qquad g_j(t)=\sum_iT_{ji}\min(u_i,t),\\
 t_0&=\min\{t\ge t_\ell:g_j(t)\ge a_j\quad\forall j\},\\
 h_i&=\min(u_i,t_0).
 \end{aligned}
\]
For nonzero rows, $g_j$ is piecewise linear between rational upper-box breakpoints. The largest required root, or $t_\ell$, determines $t_0$; only lower moments force it upward. Nonemptiness gives $T_j\ell\le b_j$. Disjoint groups independently choose mass in $[\max(a_j,T_j\ell),\min(b_j,T_jh)]$, so $\C\cap[\ell,h]$ is nonempty. Every point in this intersection has maximum $t_0$ by minimality; the entire box need not be feasible.

\noindent\textit{Minimum-cap sufficiency.}
For feasible $m$ with $t=\max_i m_i$, clipping to $v_i=\min(m_i,t_0)$ preserves lower boxes and upper moments. Restore violated lower moments within $[v,h]$ to $a_j\le b_j$. The cap permits these repairs; disjoint supports prevent interference, giving feasible $v'\ge v$. Since $p$ sums to one, $p^\top(m-v)\le t-t_0$, hence
\[
 \max_i v_i'-p^\top v'
 \le t_0-p^\top v
 \le t-p^\top m.
\]
Thus maximizing $p^\top m$ within this cap minimizes the original gap.

\textbf{Step 2: Allocate means within the cap.}
With the largest mean fixed at $t_0$, maximizing reference payoff $p^\top m$ minimizes the family upper bound. These are mean values, not policy probabilities. Increasing $m_i$ by $\epsilon$ costs group mass $T_{ji}\epsilon$ and gains $p_i\epsilon$. Starting at $\ell$, fill each nonzero row in decreasing $p_i/T_{ji}$ order on $T_{ji}>0$, to total mass $\min(b_j,T_jh)$; uncovered coordinates take $h_i$. Thus the added mass is $\min(b_j,T_jh)-T_j\ell$. Nonnegative weights justify filling even zero-ratio coordinates, and the total meets the lower bound.

The upper price $\lambda_j$ is zero when $b_j\ge T_jh$, including equality. Otherwise, use the largest support ratio if $b_j=T_j\ell$, or the last ratio receiving positive mass, including exact boundary exhaustion. This covers fixed coordinates and zero probabilities. The residual $\xi=p-T^\top\lambda$ is reference payoff after these group charges. Allocation satisfies
\[
 \begin{aligned}
 \xi_i>0&\Rightarrow m_i^*=h_i,&
 \xi_i<0&\Rightarrow m_i^*=\ell_i,\\
 \lambda_j(b_j-T_jm^*)&=0.
 \end{aligned}
\]
\textbf{Step 3: Recover the policy and original-row prices.}
Steps 1--2 give a family upper bound, not yet a robust policy. B ties for best at $m^*=(.6,.7,.7,.7)$, yet the allowed vector $(.6,.6,.7,.8)$ gives pure B a $-.075$ gain. Step 3 must match the family bound with an original-constraint policy lower bound.
We need a probability vector supported on maximizers of $m^*$ and prices
making the original-box bound~\eqref{eq:dual} tight there. The initial policy
cancels positive residuals at capped coordinates; a tight original constraint
supplies missing probability mass. Initialize
\[
 \widetilde\pi_i=(\xi_i)_+\mathbf1\{u_i\ge t_0\},\qquad
 D_\pi=1-\sum_i\widetilde\pi_i\ge0.
\]
Here $0\le\widetilde\pi_i\le p_i$, with positive mass only at maximum coordinates. The original-box coefficient $z_i^{\rm init}=\widetilde\pi_i-\xi_i$ is positive only at $\ell_i$ and negative only at $u_i$: below the cap $h_i=u_i$; elsewhere policy mass cancels positive residuals.

If $t_0=t_\ell$, add $D_\pi$ to a coordinate with $\ell_i=t_0$ and set all
lower prices $\eta$ to zero; its box coefficient increases at this lower endpoint.
Otherwise choose the first row attaining its lower bound at $t_0$.
Since $T_jh=a_j$, its supported means must equal $h_i$, and zero-price
priority gives $\lambda_j=0$, even if its upper moment is also tight. Set
\[
 E=\{i:T_{ji}>0,\ u_i\ge t_0\},\qquad
 s=\sum_{i\in E}T_{ji}>0,
\]
where positivity follows because otherwise this row would attain its bound
below $t_0$. Set $\eta_j=D_\pi/s$, with other lower prices zero.
Add $\eta_jT_{ji}$ to
$\widetilde\pi_i$ on $E$. Updated box coefficients are zero on $E$ and
$-p_i-\eta_jT_{ji}\le0$ elsewhere in this row, where $m_i^*=h_i=u_i$;
other coefficients are unchanged. The non-strict $u_i\ge t_0$ includes
roots at breakpoints.

The completed $\pi$ sums to one and selects maximum coordinates.
With $z=\pi-p+T^\top(\lambda-\eta)$, box signs and tight priced rows
make~\eqref{eq:dual} exact for $c=\pi-p$ and
$(\eta_L,\eta_U)=(\eta,\lambda)$:
\[
 \begin{aligned}
 D(\pi-p;(\eta,\lambda))&=(\pi-p)^\top m^*\\
                      &=t_0-p^\top m^*.
 \end{aligned}
\]
Scale each context's objective and both row-price types by the same $w_x^*$.
Concatenation yields original-program lower $=F(m^*)=\mathcal R(\C)$,
including zero-weight contexts. Such contexts still require a feasible point
and a legal policy; only their objective and multipliers vanish.

The receiver checks original constraints, legal policy probabilities,
multiplier signs and lower/objective/upper equality against the caller's request.
It trusts neither the artificial cap nor the producer's internal state.\hfill$\square$

\textbf{The running example, step by step.}
Using Fig.~\ref{fig:example}'s decimal inputs, Step 1 gives $t_0=.7$.
Step 2 gives $m^*=(.6,.7,.7,.7)$ and $\lambda=(0,0)$.
Step 3 starts with $(0,.25,.25,.25)$, adds the missing $.25$ to C, and
returns $\pi^*=(0,.25,.5,.25)$ with $\eta_L=(.5,0)$ and $\eta_U=0$.
The lower bound $.65(.5)-.6(.5)=.025$ equals the family upper bound
$.7-.675=.025$. Pure C in Fig.~\ref{fig:example} is another optimal policy.
This exact decimal illustration explains the construction; measured runs
check their own represented inputs.
